% Created 2026-06-21 dom 20:42
% Intended LaTeX compiler: pdflatex
\documentclass{article}
\usepackage[utf8]{inputenc}
\usepackage[T1]{fontenc}
\usepackage{graphicx}
\usepackage{longtable}
\usepackage{wrapfig}
\usepackage{rotating}
\usepackage[normalem]{ulem}
\usepackage{amsmath}
\usepackage{amssymb}
\usepackage{capt-of}
\usepackage{hyperref}

\usepackage{fancyhdr}
\usepackage[top=25mm, left=25mm, right=25mm]{geometry}
\usepackage{longtable}
\fancyhead[R]{}
\setlength\headheight{43.0pt}
\author{Juan Auz, Quillupangui Ana María, Palomeque Matías, Lanchimba Danny}
\date{2026-06-19}
\title{Taller de Comandos Linux}
\hypersetup{
 pdfauthor={Juan Auz, Quillupangui Ana María, Palomeque Matías, Lanchimba Danny},
 pdftitle={Taller de Comandos Linux},
 pdfkeywords={},
 pdfsubject={},
 pdfcreator={Emacs 27.1 (Org mode 9.7.5)}, 
 pdflang={Español}}
\begin{document}

\maketitle
\fancyhead[C]{\includegraphics[scale=0.05]{../images/logoEPN.jpg}\\
ESCUELA POLITÉCNICA NACIONAL\\FACULTAD DE INGENIERÍA DE SISTEMAS\\
ARQUITECTURA DE COMPUTADORES}
\thispagestyle{fancy}
\section{Objetivos}
\label{sec:org6dbe163}

\begin{itemize}
\item Conocer el funcionamiento del terminal de \emph{Linux}
\item Aprender a manejar comandos del terminal de \emph{Linux}
\end{itemize}
\section{Instrucciones}
\label{sec:org5d58db0}
Este taller/tutorial de comandos de Linux tiene por objetivo ejercitar
en la aplicación de comandos estudiados en clase. Para esto siga las
Instrucciones indicadas en cada uno de los numerales. Ejecute los
ejercicios llenando los comandos adecuados dentro de los bloques de
código de Emacs Org. Los resultados de la ejecución se obtienen
haciendo \texttt{C-c C-c}

Para todos los ejercicios se sugiere usar directamente un
distro de \emph{Linux} o el \emph{WSL} configurado para la clase de
\emph{Arquitectura de Computadores}. Ejecute los comandos en la terminal de
Linux. Copie los comandos usados a este archivo.

Si se ha utilizado IA para consultar temas referentes al desarrollo de
código en el punto \ref{sec:orgc7b3b1c}, adjunte el archivo \texttt{.md} con el
reporte de los diferentes prompts y respuestas dadas por la IA.
\subsection{Ejecución desde Emacs}
\label{sec:orgc31157b}

Si desea ejecutar los comandos desde Emacs debe observar que los
comandos del \texttt{shell} de Linux se ejecutarán desde la ubicación del
archivo \texttt{.org}. Por tanto, controle la ubicación de \texttt{home} en los
bloques de código. Para ejecutar los bloques de código ejecute: \texttt{C-c C-c}. 
\section{Comandos de Linux}
\label{sec:org7d800b5}
Para este taller descargaremos los datos del \emph{Fashion-Mnist}, un
dataset popular para la prueba de algoritmos de machine learning. Al
final, el dataset quedará organizado para ser aprovechado para
entrenar un modelo.

\begin{enumerate}
\item Utilice \texttt{cd} para localizarse en la posición de home de su usuario y
verifique ejecutando el comando \texttt{pwd}

\begin{verbatim}
     cd
     pwd
\end{verbatim}

\phantomsection
\label{org6b4ae11}
\begin{verbatim}
/home/juand
\end{verbatim}

\item Cree tres directorios denominados FashionData, Train y
Test. Verifique usando el comando \texttt{ls}
\begin{verbatim}
    cd ..
    mkdir FashionData
    mkdir Train
    mkdir Test
    ls
\end{verbatim}

\begin{table}[htbp]
\label{tab:org4a6f48f}
\centering
\begin{tabular}{l}
FashionData\\
FormatoTareas\\
images\\
Proyecto\\
Talleres\\
Test\\
Train\\
\end{tabular}
\end{table}

\item Investigue sobre el comando \texttt{wget} y utilícelo para descargar en
\textasciitilde{}/FashionData los datos de train y test del /Fashion-Mnist desde el
siguiente enlace: \href{https://github.com/zalandoresearch/fashion-mnist}{fashion-mnist}
\begin{enumerate}
\item Imágenes de entrenamiento:

\url{http://fashion-mnist.s3-website.eu-central-1.amazonaws.com/train-images-idx3-ubyte.gz}
\item Etiquetas de entrenamiento:

\url{http://fashion-mnist.s3-website.eu-central-1.amazonaws.com/train-labels-idx1-ubyte.gz}
\item Imágenes de prueba:

\url{http://fashion-mnist.s3-website.eu-central-1.amazonaws.com/t10k-images-idx3-ubyte.gz}
\item Etiquetas de prueba:

\url{http://fashion-mnist.s3-website.eu-central-1.amazonaws.com/t10k-labels-idx1-ubyte.gz}
\end{enumerate}
\begin{verbatim}
   cd ..
   cd FashionData
   wget http://fashion-mnist.s3-website.eu-central-1.amazonaws.com/train-images-idx3-ubyte.gz
   wget http://fashion-mnist.s3-website.eu-central-1.amazonaws.com/train-labels-idx1-ubyte.gz
   wget http://fashion-mnist.s3-website.eu-central-1.amazonaws.com/t10k-images-idx3-ubyte.gz
   wget  http://fashion-mnist.s3-website.eu-central-1.amazonaws.com/t10k-labels-idx1-ubyte.gz
   pwd
   ls 
\end{verbatim}

\begin{table}[htbp]
\label{tab:org0de98f2}
\centering
\begin{tabular}{l}
/home/juand/Documentos/ArquitecturaComp/ArquitecturaComputadores/FashionData\\
t10k-images-idx3-ubyte.gz\\
t10k-labels-idx1-ubyte.gz\\
Test\\
Train\\
train-images-idx3-ubyte.gz\\
train-labels-idx1-ubyte.gz\\
\end{tabular}
\end{table}

\item Utilice el comando \texttt{mv} de manera recursiva para pasar los datos de
entrenamiento a la carpeta \textasciitilde{}/Train
\begin{verbatim}
    cd .. 
    mv FashionData/train*  Train/
    cd Train
    ls
    

\end{verbatim}

\begin{table}[htbp]
\label{tab:org407d623}
\centering
\begin{tabular}{l}
train-images-idx3-ubyte\\
train-images-idx3-ubyte.gz\\
train-labels-idx1-ubyte\\
train-labels-idx1-ubyte.gz\\
\end{tabular}
\end{table}

\item Utilice el comando \texttt{mv} de manera recursiva para pasar los datos de
prueba a la carpeta \textasciitilde{}/Test
\begin{verbatim}
   cd ..
   mv FashionData/t10k* Test/
   cd Test
   ls

\end{verbatim}

\begin{table}[htbp]
\label{tab:orgbd0b6e0}
\centering
\begin{tabular}{l}
t10k-images-idx3-ubyte\\
t10k-images-idx3-ubyte.gz\\
t10k-labels-idx1-ubyte\\
t10k-labels-idx1-ubyte.gz\\
\end{tabular}
\end{table}

\item Verifique que la carpeta FashionData está vacía
\begin{verbatim}
    cd ..
    cd FashionData
    ls 
\end{verbatim}

\begin{table}[htbp]
\label{tab:org85ac1be}
\centering
\begin{tabular}{l}
Test\\
Train\\
\end{tabular}
\end{table}

\item Los archivos anteriores están comprimidos. Investigue como
descomprimir usando el comando \texttt{gunzip}. Si es necesario realice la
instalación del comando. Descomprima los archivos en sus
respectivas carpetas. Apunte el comando utilizado para la descompresión.
\begin{verbatim}
   cd ..
   gunzip -r Test/
   gunzip -r Train/
   cd Test/
   ls
   cd ..
   cd Train/
   ls

\end{verbatim}

\begin{table}[htbp]
\label{tab:org1d18902}
\centering
\begin{tabular}{l}
t10k-images-idx3-ubyte\\
t10k-images-idx3-ubyte.gz\\
t10k-labels-idx1-ubyte\\
t10k-labels-idx1-ubyte.gz\\
train-images-idx3-ubyte\\
train-images-idx3-ubyte.gz\\
train-labels-idx1-ubyte\\
train-labels-idx1-ubyte.gz\\
\end{tabular}
\end{table}

\item Mueva recursivamente las carpetas \textasciitilde{}/Train y \textasciitilde{}/Test dentro de \textasciitilde{}/FashionData
\begin{verbatim}
   cd ..
   mv Test/ FashionData/
   mv Train/ FashionData/
   cd FashionData/
   pwd
   ls
\end{verbatim}

\begin{table}[htbp]
\label{tab:org85d011f}
\centering
\begin{tabular}{l}
/home/juand/Documentos/ArquitecturaComp/ArquitecturaComputadores/FashionData\\
Test\\
Train\\
\end{tabular}
\end{table}

\item Verifique la estructura de la carpeta \texttt{/FashionData usando el
   comando \textasciitilde{}tree}
\begin{verbatim}
   cd ..
   cd FashionData
   tree
\end{verbatim}

\begin{table}[htbp]
\label{tab:orgca4e6cf}
\centering
\begin{tabular}{llll}
. &  &  & \\
├── & Test &  & \\
│   & ├── & t10k-images-idx3-ubyte & \\
│   & └── & t10k-labels-idx1-ubyte & \\
└── & Train &  & \\
├── & train-images-idx3-ubyte &  & \\
└── & train-labels-idx1-ubyte &  & \\
 &  &  & \\
3 & directories, & 4 & files\\
\end{tabular}
\end{table}
\end{enumerate}
\section{Problema}
\label{sec:orgc7b3b1c}
1
Se desea realizar un registro climatológico de la ciudad de
Quito. Para esto, escriba un script de Python/Java que permita obtener datos climatológicos desde el API de
\href{https://openweathermap.org/current\#one}{openweathermap}. Considere la latitud y la longitud de Quito como
(-0.2299, -78.5249), respectivamente. Los resultados obtenidos de
la consulta al API se escriben en un archivo
\emph{clima-quito-hoy.csv}. Cada ejecución del script debe almacenar
nuevos datos en el archivo. Investigue sobre \textbf{crontab} y utilícelo
para obtener datos del API de \emph{openweathermap} cada 5 minutos
mediante la ejecución de un archivo ejecutable denominado
\emph{get-weather.sh}. Verifique los resultados. Todas las operaciones
se realizan en Linux o en el WSL. Las etapas del problema se
subdividen en:

\begin{enumerate}
\item Crear su API gratuito en \href{https://openweathermap.org/current\#one}{openweathermap}
\item Escribir un script en Python/Java que realice la consulta al API y
escriba los resultados en \emph{clima-quito-hoy.csv}

\begin{verbatim}
import requests
import json
import csv
import os
import pathlib
from datetime import datetime

API_KEY = 'b794b9218f841b4d649bce12eacc5aed'
LAT = -0.2299
LON = -78.5249
URL = f"https://api.openweathermap.org/data/2.5/weather?lat={LAT}&lon={LON}&appid={API_KEY}&units=metric"
OUTPUT_FILE = pathlib.Path(__file__).parent / 'clima-quito-hoy.csv'

def main():
    response = requests.get(URL)
    data = response.json()
    timestamp = datetime.now().isoformat()

    existe = os.path.isfile(OUTPUT_FILE)

    with open(OUTPUT_FILE, 'a', newline='', encoding='utf-8') as f:
        writer = csv.writer(f)
        if not existe:
            writer.writerow(['timestamp', 'data'])
        writer.writerow([timestamp, json.dumps(data)])

    print(f"[{timestamp}] Guardado - temp: {data['main']['temp']}°C")

if __name__ == '__main__':
    main()

\end{verbatim}

\phantomsection
\label{org0aa7755}
\begin{verbatim}
None
\end{verbatim}

\item Desarrollar un ejecutable \emph{get-weather.sh} para ejecutar el
programa Python/Java.

\begin{verbatim}
   sudo nano get-weather.sh
   
   #!/bin/bash
   /home/juand/miniforge3/envs/iccd332/bin/python /home/juand/Documentos/ArquitecturaComp/ArquitecturaComputadores/Proyecto/get-weather.py
   
   chmod +x get-weather.sh
   ./get-weather.sh
\end{verbatim}

\phantomsection
\label{orgbdd3f96}
\begin{verbatim}
[2026-06-21T20:42:30.553918] Guardado - temp: 10.64°C
\end{verbatim}

\item Configurar Crontab para la adquisición de datos. Escriba el comando
configurado.
*/5 * * * * /home/juand/Documentos/ArquitecturaComp/ArquitecturaComputadores/Proyecto/get-weather.sh >> /home/juand/Documentos/ArquitecturaComp/ArquitecturaComputadores/Proyecto/cron.log 2>\&1
\end{enumerate}
\end{document}
